% ===================================================================
\section{Natural selection (EE10)}
% ===================================================================

\subsection{Definition and Darwin's four postulates}
\textbf{Natural selection is the differential survival and reproduction of individuals due to differences in
phenotype} (Zimmer and Emlen). The natural/artificial distinction ``is rather\dots artificial'' and persists for
historical reasons: it is one process. Freeman and Herron's form of Darwin's postulates:
\begin{enumerate}
\item individuals in a population \textbf{vary};
\item some of the variation is \textbf{passed from parents to offspring};
\item some individuals are \textbf{more successful} at surviving and reproducing;
\item success is \textbf{not mere luck}: it depends on the inherited variant traits.
\end{enumerate}
If all four hold, the composition of the population \emph{inevitably} changes across generations.

\begin{keyidea}[: selection needs no ``struggle'' in the literal sense]
The postulates require differences in \emph{reproductive success}. Genotypes can differ in fecundity, mating success or
development speed with everyone surviving and food abundant; relative output decides. Darwin used ``struggle'' in a
large, metaphorical sense (2023 end-sem Q2).
\end{keyidea}

\subsection{The medium ground finch on Daphne Major}
\emph{Geospiza fortis}, monitored yearly since \textbf{1973} by Peter and Rosemary Grant ($\sim$1200 birds, most marked by 1977).

\begin{center}\small
\begin{tabularx}{\textwidth}{lX}
\toprule
postulate & evidence\\\midrule
1 variation & beak depth varies, roughly bell-shaped ($\sim$7--11\,mm)\\
2 heritability & deep-beaked parents have deep-beaked chicks; after DNA showed 20\,\% extra-pair chicks and these were excluded,
$h^2=0.65\pm0.15$\\
3 differential survival & \textbf{1977 drought}: 24\,mm rain instead of $\sim$130\,mm; no breeding; \textbf{84\,\%} of \emph{G.\ fortis}
disappeared, mostly starving\\
4 non-random & small soft seeds ran out first; the hard fruits of \emph{Tribulus cistoides} remained, crackable only by large,
deep-beaked birds; survivors had deeper beaks and bigger bodies\\
consequence & 1978 chicks had deeper beaks than 1976 chicks: \textbf{the population evolved}\\\bottomrule
\end{tabularx}
\end{center}

Then the direction reversed: the \textbf{1983 El Niño} (1359\,mm, $\sim$57$\times$ 1977's rain) made small soft seeds
abundant, favouring smaller birds. In the \textbf{2003--04 drought}, the newly arrived large ground finch \emph{G.\
magnirostris} took the big \emph{Tribulus} seeds, so \emph{G.\ fortis} evolved \emph{smaller} beaks: competition changed the
direction of selection.

\begin{chain}[: why extra-pair paternity mattered]
A chick sired by an extra-pair male is plotted against a social father it is unrelated to. That adds noise to the
parent--offspring regression and flattens it, so misassigned paternity \emph{underestimates} heritability. Removing those
chicks gave the higher, truer $h^2=0.65$.
\end{chain}

\textbf{Three lessons} (slide 3): (1) evolution can be \textbf{rapid and dramatic}; (2) sources of selection
\textbf{shift}, so evolution reverses or changes direction; (3) the \textbf{see-saw} in a changing environment can look like
stately gradual change in fossils: a fossil bed mixes thousands of years, the zig-zags largely cancel, and only the small
net change remains visible. Slow change in the record is compatible with fast change each generation.

\subsection{Three modes of selection}
\begin{center}
\begin{tikzpicture}
\begin{groupplot}[group style={group size=3 by 1,horizontal sep=8mm},width=5.4cm,height=3.6cm,
  ticks=none,xmin=-4,xmax=4,ymin=0,ymax=0.62,axis lines=left]
\nextgroupplot[title={\small directional}]
\addplot[gray,thick,domain=-4:4,samples=80]{exp(-x^2/2)/sqrt(2*pi)};
\addplot[NavyBlue,thick,domain=-4:4,samples=80]{exp(-(x-1)^2/2)/sqrt(2*pi)};
\nextgroupplot[title={\small stabilizing}]
\addplot[gray,thick,domain=-4:4,samples=80]{exp(-x^2/2)/sqrt(2*pi)};
\addplot[NavyBlue,thick,domain=-4:4,samples=80]{exp(-x^2/(2*0.5))/sqrt(2*pi*0.5)};
\nextgroupplot[title={\small disruptive}]
\addplot[gray,thick,domain=-4:4,samples=80]{exp(-x^2/2)/sqrt(2*pi)};
\addplot[NavyBlue,thick,domain=-4:4,samples=120]{0.5*exp(-(x-1.4)^2/(2*0.5))/sqrt(2*pi*0.5)+0.5*exp(-(x+1.4)^2/(2*0.5))/sqrt(2*pi*0.5)};
\end{groupplot}
\end{tikzpicture}

\small Grey: trait distribution before; blue: after. Directional: the mean moves. Stabilizing: intermediates favoured,
variance falls. Disruptive: both extremes favoured, variance rises, may split.
\end{center}

\subsection{The nature of selection}
\begin{chain}[: selection on individuals, evolution of populations]
\begin{enumerate}
\item An individual's genotype is fixed for life. It survives or dies, reproduces more or less, but does not change.
No finch's beak grew from cracking \emph{Tribulus}.
\item Selection therefore \emph{sorts} individuals; the \emph{population's} composition shifts as a result.
\item Selection acts on \textbf{phenotypes}, but evolution happens only if \textbf{allele frequencies} change, which
requires the phenotypic differences to be heritable. With $h^2=0$ the drought would still have killed the small-beaked
birds, but the next generation would have bounced back: no evolution.
\item So ``selection acts on phenotypes; genotypes get selected; genes get differentially transmitted'' (wrap-up slide 13),
and \textbf{natural selection is not evolution}: selection without heritability gives no evolution, and evolution also has
other causes (drift, mutation, migration, non-random mating: the broken HW assumptions).
\end{enumerate}
\end{chain}
``The mess that is natural selection'' (slide 10): genes $\to$ genotype $\to$ (with biotic and abiotic environment)
phenotype $\to$ selection $\to$ meiosis $\to$ fertilisation $\to$ genes, with \textbf{chance} at every step.

\begin{trap}[: heritability is not ``resemblance'']
``Obese children often have obese parents'' or ``long-winged birds have long-winged offspring'' does not show a genetic
basis: shared diet, island, rearing conditions or maternal effects produce resemblance too. A genetic basis needs a
common-garden or cross-fostering design, or crosses. And predicting a response needs differential fitness as well
(2022 end-sem Q4; 2024 end-sem Q2).
\end{trap}

\subsection{What selection is not}
\begin{itemize}
\item \textbf{Always behind}: it changes frequencies \emph{after} exposure; the 1978 finches were adapted to 1977, not to
what came next. It cannot anticipate; new mutations do not arise to meet needs; a population facing too-fast change may go
extinct.
\item \textbf{Not progressive}: organisms ``devolve'' as often as they evolve (parasites); adaptation proceeds by gain,
loss and modification of characters.
\item \textbf{Works with what is available}: pre-existing variation, co-opted structures (\textbf{preadaptations}); no grand
design, no perfection. \emph{Gambusia} mosquitofish: a large gonopodium attracts mates but hinders escape, so males in
high-predation populations have smaller ones: a \textbf{trade-off}, not an optimum.
\item \textbf{Yet new traits evolve}: mutation and recombination create new combinations, which selection sorts (Illinois
corn oil: from 4--6\,\% to $\sim$20\,\% in 100 generations, far beyond the starting range).
\end{itemize}

\subsection{One-locus selection (beyond the slides, but explains the past-paper items)}
Selection enters the HW cycle between zygote and adult: zygotes form at $p^2:2pq:q^2$, then survive with fitnesses
$w_{11},w_{12},w_{22}$:
\[
p'=\frac{p^2w_{11}+pq\,w_{12}}{\bar w},\qquad \bar w=p^2w_{11}+2pq\,w_{12}+q^2w_{22},\qquad
\Delta p=\frac{pq\,[\,p(w_{11}-w_{12})+q(w_{12}-w_{22})\,]}{\bar w}.
\]
\begin{chain}[: what the formula says]
\begin{itemize}
\item The factor $pq$: selection is slow when either allele is rare (little variation to act on).
\item A \textbf{recessive} favoured allele is invisible when rare (almost all copies sit in heterozygotes, HW property 5),
so it takes very long to start spreading. A \textbf{dominant} favoured allele spreads fast at first but stalls at the end,
because the last disfavoured recessive copies hide in heterozygotes. Additive alleles are fastest overall ($s=0.1$: 175
generations vs $\sim$1000).
\item \textbf{Overdominance} (heterozygote advantage, $w=1-s:1:1-t$): stable polymorphism at $p^*=t/(s+t)$; variation is
maintained (classic case: sickle-cell in malarial regions).
\item \textbf{Underdominance} (heterozygote disadvantage): an \emph{unstable} interior equilibrium; whichever allele starts
above it goes to fixation (a threshold, like the Allee effect).
\item \textbf{Breeder's equation} $R=h^2S$: only the heritable fraction of the selection differential is passed on.
\end{itemize}
\end{chain}

\begin{statements}
\begin{enumerate}[label=\arabic*.]
\item ``The finches' beaks deepened because birds needed to crack hard seeds.''
\F No individual changed; small-beaked birds died, and deep beaks were heritable.
\item ``Natural selection always requires mortality.''
\F Differences in reproduction alone suffice.
\item ``If a trait is under strong selection but has zero heritability, the population will still evolve, only more slowly.''
\F No evolution at all; the next generation returns to the old distribution.
\item ``Stabilizing selection shifts the mean toward the optimum and widens the distribution.''
\F It keeps the mean and narrows the distribution; disruptive selection widens it.
\item ``The reversal of selection in 1983 shows that the 1977 evolution was an error of measurement.''
\F Selection tracks a changing environment; reversals are expected (lesson 2).
\item ``Fossil sequences showing slow change mean evolution is slow.''
\F Fast zig-zags average out over the long spans a fossil bed mixes (lesson 3).
\item ``A new beneficial recessive allele spreads quickly because it is beneficial.''
\F It is almost never expressed while rare (it sits in heterozygotes), so selection barely sees it (2024 end-sem Q11).
\item ``Heterozygote advantage maintains both alleles in the population.''
\T Stable equilibrium at $p^*=t/(s+t)$.
\item ``Selection produces the variation it acts on.''
\F Mutation and recombination do; selection sorts.
\item ``Competition from \emph{G.\ magnirostris} changed the direction of selection on \emph{G.\ fortis} beaks.''
\T In 2004 it took the large seeds, so smaller beaks were favoured.
\end{enumerate}
\end{statements}

\begin{longans}[: ``Natural selection is not evolution. Discuss.'']
(1) Definitions: selection (differential survival/reproduction due to phenotype) vs evolution (change in allele/genotype
frequencies across generations). (2) Four postulates; Daphne Major tests each (numbers). (3) Selection acts on
individuals and phenotypes; evolution requires heritability ($h^2=0$ thought experiment; $R=h^2S$). (4) Other causes of
evolution (drift, mutation, migration, non-random mating). (5) Properties: always behind, not progressive, uses what is
available (\emph{Gambusia}), yet new traits arise. (6) Lessons: rapid, reversible, zig-zags look gradual.
\end{longans}
